\section{Scaling the Base Model: Qwen2.5-3B vs Qwen3-30B-A3B-Base}\label{app:qwen3-comparison}

To test whether a stronger base model improves the metric on Tevet's benchmarks, we re-ran the full evaluation using Qwen3-30B-A3B-Base (a 30B-parameter mixture-of-experts model with $\sim$3B active parameters per token) via the Tinker API, with the same setup as in Section~\ref{sec:tevet}: completion-format prompting, 50 permutations, no fine-tuning.
Table~\ref{tab:qwen3-comparison} reports per-task $C \times a_n$ (per-byte) for both base models.

\begin{table}[h]
\centering
\caption{Per-task comparison of $C \times a_n$ (per-byte) for Qwen2.5-3B vs Qwen3-30B-A3B-Base on Tevet diversity-eval. $\Delta$AUC is positive when the larger model wins.
ROC AUC is reported for the binary tasks (McDiv, McDiv\_nuggets, ConTest); only Spearman $\rho$ is meaningful for DecTest (continuous temperature labels).}
\label{tab:qwen3-comparison}
\small
% Generated by: scripts/compare_qwen_models.py --emit-latex
% Data sources: figures/tevet_validation/c_ainf_analysis_v3/summary_table.txt
%               figures/tevet_validation/c_ainf_analysis_qwen3/summary_table.txt
% Metric: C x a_n (per-byte). Q2.5-3B = Qwen2.5-3B local; Q3-30B = Qwen3-30B-A3B-Base via Tinker.
\begin{tabular}{@{}llrrrrr@{}}
\toprule
 &  & \multicolumn{2}{c}{Qwen2.5-3B} & \multicolumn{2}{c}{Qwen3-30B-A3B} & \\
\cmidrule(lr){3-4} \cmidrule(lr){5-6}
Section & Task & $\rho$ & AUC & $\rho$ & AUC & $\Delta$AUC \\
\midrule
McDiv & prompt\_gen (no\_hds) & 0.729 & 0.921 & 0.718 & 0.915 & -0.006 \\
 & resp\_gen (no\_hds) & 0.500 & 0.788 & 0.483 & 0.779 & -0.009 \\
 & story\_gen (no\_hds) & 0.523 & 0.802 & 0.493 & 0.785 & -0.017 \\
\midrule
McDiv\_nuggets & prompt\_gen (no\_hds) & 0.636 & 0.867 & 0.615 & 0.855 & -0.012 \\
 & prompt\_gen (with\_hds) & 0.683 & 0.895 & 0.655 & 0.878 & -0.017 \\
 & resp\_gen (no\_hds) & 0.345 & 0.699 & 0.325 & 0.687 & -0.012 \\
 & resp\_gen (with\_hds) & 0.437 & 0.753 & 0.425 & 0.746 & -0.007 \\
 & story\_gen (no\_hds) & 0.317 & 0.683 & 0.293 & 0.669 & -0.014 \\
 & story\_gen (with\_hds) & 0.405 & 0.734 & 0.376 & 0.717 & -0.017 \\
\midrule
ConTest & prompt\_gen (with\_hds) & 0.584 & 0.837 & 0.592 & 0.842 & +0.005 \\
 & resp\_gen (with\_hds) & 0.391 & 0.726 & 0.333 & 0.692 & -0.034 \\
 & story\_gen (with\_hds) & 0.686 & 0.896 & 0.653 & 0.877 & -0.019 \\
\midrule
DecTest & prompt\_gen (no\_hds) & 0.842 & --- & 0.877 & --- & --- \\
 & prompt\_gen (with\_hds) & 0.845 & --- & 0.875 & --- & --- \\
 & resp\_gen (no\_hds) & 0.771 & --- & 0.804 & --- & --- \\
 & resp\_gen (with\_hds) & 0.760 & --- & 0.777 & --- & --- \\
 & story\_gen (no\_hds) & 0.763 & --- & 0.771 & --- & --- \\
 & story\_gen (with\_hds) & 0.785 & --- & 0.785 & --- & --- \\
\bottomrule
\end{tabular}

\end{table}

\paragraph{Result: scaling did not help on the binary tasks.} On the \qwenThreeBinaryTotal{} binary classification tasks (McDiv, McDiv\_nuggets, ConTest), Qwen2.5-3B wins on \qwenThreeBinaryWinsQwenTwoFive{} of \qwenThreeBinaryTotal{}, with ConTest prompt\_gen the only marginal win for Qwen3-30B-A3B ($\qwenThreeConTestPromptGenDeltaAUC$ AUC).
The mean degradation from scaling is small but consistent (mean $\Delta$AUC = $\qwenThreeBinaryMeanDeltaAUC$ across the \qwenThreeBinaryTotal{} binary tasks).
On the \qwenThreeDecTestTotal{} DecTest tasks (continuous temperature labels), Qwen3-30B-A3B is mildly better on \qwenThreeDecTestWins{} of \qwenThreeDecTestTotal{} (mean $\qwenThreeDecTestMeanDeltaRho$ in Spearman $\rho$, with one tie).

\paragraph{Interpretation.} A natural reading is that the binary McDiv/ConTest labels have low information content: they distinguish only ``high'' vs ``low'' diversity, and this distinction is coarse enough that a 3B base model already captures the relevant signal.
A larger model has nothing more to extract that the binary label can reward.
The continuous DecTest labels are noise-free (temperature is the ground truth), so additional metric resolution from a larger model translates into measurable correlation gains.
If this interpretation is correct, the binary tasks are saturated for our metric and the path to improvement runs through finer-grained labels rather than larger base models.

\paragraph{Caveat: prompt structure was not optimized.} We made no attempt to optimize the prompt structure used to elicit ICL on these benchmarks; the formatting follows the neutral template from Section~\ref{sec:practical} verbatim.
A targeted prompt structure (e.g., explicitly framing the responses as ``Five high-diversity continuations,'' or providing in-context exemplars of diverse vs.\ repetitive sets) would likely improve the metric's discriminative power, particularly for the larger model whose ICL capabilities benefit more from carefully constructed contexts.
We leave prompt-structure optimization to future work.

